\appendix
\part{Anhang}

\newcommand{\lit}[3]{\vspace{5pt}\noindent\textbf{[#1]} #2\par\nopagebreak\vspace{1pt}{\small\itshape\color{grau} #3}\par\vspace{2pt}}

\chapter{Formelsammlung und Übersicht}

\section{Grundbegriffe}

\begin{merksatz}
\boxtitel{Spiele und Gleichgewichte}
\begin{center}
{\small
\renewcommand{\arraystretch}{1.62}
\begin{tabular}{@{}P{4.4cm}P{4.2cm}@{}}
\toprule
Spieler, $\Str_i$, $\Aus_i$ & Normalform (1.1)\\
$\Aus_1+\Aus_2 = 0$ & Nullsummenspiel (1.2)\\
$\max_r\min_c \le \min_c\max_r$ & Satz 1.3, stets\\
Zeilenmin $=$ Spaltenmax & Sattelpunkt (1.4)\\
strikt besser gegen alles & Dominanz (1.6)\\
$\vek p \ge 0$, $\sum p_i = 1$ & gemischt (2.1)\\
$\vek p^{\!\top}A\vek q$ & erwartete Auszahlung\\
\bottomrule
\end{tabular}}
\end{center}
\end{merksatz}

\begin{merksatz}
\boxtitel{Nullsummenspiele}
\begin{center}
{\small
\renewcommand{\arraystretch}{1.62}
\begin{tabular}{@{}P{4.4cm}P{4.2cm}@{}}
\toprule
$\max_{\vek p}\min_{\vek q} = \min_{\vek q}\max_{\vek p}$ & Minimax (2.2)\\
$N = a+d-b-c$ & Nenner\\
$p = \dfrac{d-c}{N}$, $q = \dfrac{d-b}{N}$ & $2\times2$ (2.3)\\
$v = \dfrac{ad-bc}{N}$ & Determinante\\
$\min \vek 1^{\!\top}\vek x$, $A^{\!\top}\vek x \ge \vek 1$ & primal (3.1)\\
$\max \vek 1^{\!\top}\vek y$, $A\vek y \le \vek 1$ & dual\\
$\vek x = \vek p/v$, $\vek y = \vek q/v$ & Rückrechnung\\
benutzt $\Rightarrow$ Ertrag $=v$ & Ausgleich (3.3)\\
\bottomrule
\end{tabular}}
\end{center}
\end{merksatz}

\begin{merksatz}
\boxtitel{Gleichgewichte}
\begin{center}
{\small
\renewcommand{\arraystretch}{1.62}
\begin{tabular}{@{}P{4.4cm}P{4.2cm}@{}}
\toprule
vier Forderungen & Nutzen (4.1)\\
bis auf $\alpha\Aus+\beta$, $\alpha>0$ & eindeutig\\
niemand gewinnt einseitig & Nash (4.3)\\
$T>R>P>S$ & Dilemma\\
$T>R>S>P$ & Feigheit\\
$R>T\ge P>S$ & Hirschjagd\\
$q^{*} = T/R$ & Vertrauensschwelle\\
$K$ konvex, kompakt, $f$ stetig & Brouwer (6.1)\\
$T(\vek\sigma)_{i,k} = \dfrac{\sigma_{i,k}+g_{i,k}}{1+\sum_j g_{i,j}}$ & Nash-Abbildung\\
$(2^{m}-1)(2^{n}-1)$ & Trägerpaare\\
\bottomrule
\end{tabular}}
\end{center}
\end{merksatz}

\section{Zeit und Evolution}

\begin{merksatz}
\boxtitel{Über die Zeit}
\begin{center}
{\small
\renewcommand{\arraystretch}{1.62}
\begin{tabular}{@{}P{4.4cm}P{4.2cm}@{}}
\toprule
vollständiger Plan & Strategie (7.1)\\
rückwärts vom Blatt & Rückwärtsinduktion\\
in jedem Teilspiel Nash & perfekt (7.3)\\
$\sum_t \delta^{t} u_t$ & Abzinsung\\
$1/(1-\delta)$ & erwartete Runden\\
$\delta \ge \dfrac{T-R}{T-P}$ & Grimm (8.2)\\
$\delta \ge \dfrac{T-R}{R}$ & Tit for Tat\\
erreichbar $+$ über Minmax & Volkssatz (8.3)\\
\bottomrule
\end{tabular}}
\end{center}
\end{merksatz}

\begin{merksatz}
\boxtitel{Evolution}
\begin{center}
{\small
\renewcommand{\arraystretch}{1.62}
\begin{tabular}{@{}P{4.4cm}P{4.2cm}@{}}
\toprule
$\Aus(\sigma^{*},\sigma^{*}) > \Aus(\sigma,\sigma^{*})$ & ESS, Fall (i)\\
oder $=$ und $\Aus(\sigma^{*},\sigma)>\Aus(\sigma,\sigma)$ & Fall (ii)\\
$p^{*} = V/C$ & Falke-Taube\\
$\dfrac{V(C-V)}{2C}$ & Eignung im GG\\
$\dot x_i = x_i(f_i-\bar f)$ & Replikator (10.1)\\
$\dot x = x(1-x)(f_1-f_2)$ & zwei Typen\\
$H = \sum p_i \ln x_i$ & erhalten bei $A^{\!\top}=-A$\\
$T_{\text{aus}} \ge \dfrac{r_{\text{an}}}{|r_{\text{aus}}|}T_{\text{an}}$ & adaptive Therapie\\
\bottomrule
\end{tabular}}
\end{center}
\end{merksatz}

\section{Regeln entwerfen}

\begin{merksatz}
\boxtitel{Kollektivgüter und Wahlen}
\begin{center}
{\small
\renewcommand{\arraystretch}{1.62}
\begin{tabular}{@{}P{4.4cm}P{4.2cm}@{}}
\toprule
$\Aus_i = e-c_i+\dfrac rn\sum_j c_j$ & Kollektivgut\\
$-1+r/n$ & Grenzertrag\\
$1<r<n$ & Dilemma\\
$X_{\text{GG}} = \dfrac{n(a-c)}{b(n+1)}$ & Allmende\\
$X^{*} = \dfrac{a-c}{2b}$ & Optimum\\
$2n/(n+1)$ & Übernutzung (11.2)\\
$\tau = b\,(X-x_i)$ & Pigou-Abgabe\\
\midrule
Mehrheiten nicht übertragbar & Condorcet\\
vier Forderungen, keine Lösung & Arrow (12.1)\\
keines manipulationsfest & Gibbard (12.2)\\
Median gewinnt & eingipflig (12.3)\\
\bottomrule
\end{tabular}}
\end{center}
\end{merksatz}

\begin{merksatz}
\boxtitel{Auktionen, Zuteilung, Verhandlung}
\begin{center}
{\small
\renewcommand{\arraystretch}{1.62}
\begin{tabular}{@{}P{4.4cm}P{4.2cm}@{}}
\toprule
Wahrheit dominant & Zweitpreis (13.1)\\
$b(v) = \dfrac{n-1}{n}v$ & Erstpreis\\
$\dfrac{n-1}{n+1}$ & Erlös, alle Formen\\
höchstes Signal $>$ Wert & Fluch des Siegers\\
Anfragen, halten, ablehnen & Gale-Shapley\\
anfragende Seite bestes & Satz 13.4\\
nie für beide Seiten & Satz 13.5\\
\midrule
$\max (u_1-d_1)(u_2-d_2)$ & Nash (14.1)\\
$x^{*} = \dfrac{1-\delta_2}{1-\delta_1\delta_2}$ & Rubinstein (14.2)\\
$\sum_{i\in S}x_i \ge v(S)$ & Kern (14.4)\\
mittlerer Beitrag über $n!$ & Shapley (14.5)\\
konvex $\Rightarrow$ Shapley im Kern & Satz von 1971\\
\bottomrule
\end{tabular}}
\end{center}
\end{merksatz}

\section{Die Landkarte des Bandes}

\begin{merksatz}
\boxtitel{Welcher Satz wo --- und worauf er beruht}
\begin{center}
{\small
\renewcommand{\arraystretch}{1.42}
\begin{tabular}{@{}P{3.4cm}P{1.2cm}P{3.6cm}@{}}
\toprule
\textbf{Ergebnis} & \textbf{Kap.} & \textbf{Beruht auf}\\
\midrule
Maximin $\le$ Minimax & $1$ & Definition von max und min\\
Sattelpunkt $\iff$ gleich & $1$ & Satz 1.3\\
\textbf{Minimax-Satz} & $2$ & Brouwer (von Neumann)\\
$2\times2$ ausrechnen & $2$ & lineare Gleichung\\
Minimax $=$ Dualität & $3$ & Band 19, Satz 4.3\\
Ausgleichssatz & $3$ & Mittelwertargument\\
Nutzenfunktion & $4$ & vier Axiome\\
\textbf{Existenz} & $6$ & Brouwer $+$ Mittelwert\\
$\mathsf{PPAD}$-Vollständigkeit & $6$ & Chen, Deng 2006\\
Zermelo & $7$ & Endlichkeit\\
Teilspielperfektheit & $7$ & Rückwärtsinduktion\\
endlich zerstört alles & $8$ & Rückwärtsinduktion\\
Grimm-Schwelle & $8$ & geometrische Reihe\\
Volkssatz & $8$ & Konstruktion\\
ESS $\Rightarrow$ Nash & $9$ & Fallunterscheidung\\
Replikator & $10$ & Quotientenregel\\
geschlossene Bahnen & $10$ & Antisymmetrie\\
Grenzertrag & $11$ & Ableitung\\
Übernutzung & $11$ & Ableitung\\
\textbf{Arrow} & $12$ & entscheidende Gruppen\\
Medianwähler & $12$ & Eingipfligkeit\\
Vickrey & $13$ & Fallunterscheidung\\
Erlösgleichheit & $13$ & Ordnungsstatistik\\
Gale-Shapley & $13$ & Widerspruch\\
Verhandlungslösung & $14$ & vier Axiome\\
Rubinstein & $14$ & zwei Gleichungen\\
\bottomrule
\end{tabular}}
\end{center}
\end{merksatz}

\begin{merksatz}
\boxtitel{Was in diesem Band \kursiv{nicht} bewiesen wurde}
\begin{center}
{\small
\renewcommand{\arraystretch}{1.5}
\begin{tabular}{@{}P{3.4cm}P{1.2cm}P{3.4cm}@{}}
\toprule
\textbf{Ergebnis} & \textbf{Kap.} & \textbf{Wo es steht}\\
\midrule
Fixpunktsatz von Brouwer & $6$ & Osborne\zit{19}\\
Satz von Arrow & $12$ & Osborne\zit{19}, Kap.\ 10\\
Nutzensatz (4.1) & $4$ & von Neumann\zit{18}\\
Minimax allgemein & $2$ & über Kapitel 3\\
Dualitätssatz & $3$ & Band 19, Kap.\ 4\\
Satz von Wilson & $4$ & Myerson\zit{16}\\
Satz von Shapley (1971) & $14$ & Peters\zit{21}\\
Gibbard-Satterthwaite & $12$ & Myerson\zit{16}\\
Erstpreis-Gleichgewicht & $13$ & Krishna\zit{13}\\
\bottomrule
\end{tabular}}
\end{center}

\begin{center}
\begin{minipage}{0.88\textwidth}
\kursiv{Neun Auslassungen, jede im Text benannt, mit Begründung und Fundstelle --- die einzige zulässige Form, die Regel dieses Werkes zu brechen.}
\end{minipage}
\end{center}
\end{merksatz}

\section{Zeittafel}

\begin{merksatz}
\boxtitel{Von einem Brief über Kartenspiele bis zur Erklärbarkeit}
\begin{center}
{\small
\renewcommand{\arraystretch}{1.34}
\begin{tabular}{@{}P{1.4cm}P{2.8cm}P{4.0cm}@{}}
\toprule
\textbf{Jahr} & \textbf{Person} & \textbf{Ereignis}\\
\midrule
$1713$ & Waldegrave & erste Minimax-Lösung, brieflich\\
$1755$ & Rousseau & die Hirschjagd\\
$1770$ & Borda & Punkteverfahren\\
$1785$ & Condorcet & das Paradoxon\\
$1833$ & Lloyd & die offene Weide\\
$1911$ & Brouwer & Fixpunktsatz\\
$1912$ & Zermelo & Satz über das Schachspiel\\
$1921$ & Borel & gemischte Strategien\\
$1928$ & \textbf{von Neumann} & \textbf{Minimax-Satz}\\
$1941$ & Kakutani & Fixpunktsatz für Mengen\\
$1944$ & von Neumann, Morgenstern & das Buch\\
$1947$ & Dantzig, von Neumann & Dualität erkannt\\
$1948$ & Black & Medianwähler\\
$1950$ & Flood, Dresher & das erste Dilemma-Experiment\\
$1950$ & \textbf{Nash} & \textbf{Gleichgewicht, Verhandlung}\\
$1950$ & Tucker & der Name\\
$1951$ & \textbf{Arrow} & \textbf{Unmöglichkeitssatz}\\
$1953$ & Kuhn & extensive Form\\
$1953$ & Shapley & der Wert\\
$1953$ & Allais & das Paradoxon\\
$1960$ & Schelling & Festlegung und Drohung\\
$1961$ & Vickrey & Zweitpreisauktion\\
$1962$ & Gale, Shapley & stabile Zuordnung\\
$1965$ & Selten & Teilspielperfektheit\\
$1965$ & Olson & Kollektivgüter\\
$1967$ & Harsanyi & unvollständiges Wissen\\
$1968$ & Hardin & Tragik der Allmende\\
$1971$ & Wilson & ungerade Anzahl\\
$1973$ & \textbf{Maynard Smith, Price} & \textbf{stabile Strategien}\\
$1973$ & Gibbard & Manipulierbarkeit\\
$1974$ & Aumann & korreliertes Gleichgewicht\\
$1978$ & Taylor, Jonker & Replikatorgleichung\\
$1980$ & Axelrod & die Turniere\\
$1982$ & Rubinstein & Verhandlungsmodell\\
$1982$ & Roth & Zuordnung und Wahrheit\\
$1990$ & Ostrom & Selbstverwaltung\\
$1992$ & Nowak, May & räumliche Spiele\\
$1994$ & --- & Frequenzauktionen\\
$1996$ & Sinervo, Lively & zyklische Eidechsen\\
$1999$ & Fehr, Schmidt & Ungleichheitsaversion\\
$2006$ & Chen, Deng & $\mathsf{PPAD}$-Vollständigkeit\\
$2009$ & Gatenby & adaptive Therapie\\
$2017$ & Lundberg, Lee & Shapley zur Erklärbarkeit\\
\bottomrule
\end{tabular}}
\end{center}
\end{merksatz}

\section{Glossar}

\begin{merksatz}
\boxtitel{Begriffe von A bis Z}
\begin{center}
{\small
\renewcommand{\arraystretch}{1.3}
\begin{longtable}{@{}P{3.6cm}P{5.0cm}@{}}
\toprule
\textbf{Begriff} & \textbf{Bedeutung}\\
\midrule
\endhead
Allmende & gemeinsam genutzte Ressource (11)\\
Auszahlung & Nutzen eines Ausgangs (1.1)\\
Beste Antwort & maximiert gegen gegebene Wahl (4.2)\\
Borda-Verfahren & Punkte nach Rangplatz (12)\\
Condorcet-Sieger & schlägt jeden im Paarvergleich (12)\\
Dominanz & besser gegen jede Wahl (1.6)\\
Eingipflig & Bestpunkt, nach beiden Seiten fallend (12)\\
Eignung & Auszahlung als Fortpflanzungserfolg (9)\\
ESS & evolutionär stabil (9.1)\\
Extensive Form & Spielbaum (7.1)\\
Fluch des Gewinners & Sieger hat überschätzt (13)\\
Gemischte Strategie & Wahrscheinlichkeitsvektor (2.1)\\
Grimm & unwiderrufliche Vergeltung (8)\\
Informationsmenge & ununterscheidbare Knoten (7)\\
Kern & keine Gruppe schert aus (14.4)\\
Kollektivgut & Beitrag nützt allen (11)\\
Konfliktpunkt & Ergebnis ohne Einigung (14)\\
Minmax-Wert & was man sich sichern kann (8)\\
Nash-Gleichgewicht & niemand gewinnt einseitig (4.3)\\
Normalform & Matrixdarstellung (1.1)\\
Nullsummenspiel & Gewinn $=$ Verlust (1.2)\\
Pareto-optimal & niemand besser ohne Verschlechterung\\
Replikatorgleichung & Wachstum nach Überdurchschnitt (10.1)\\
Risikodominanz & sicherer trotz geringerem Ertrag (5)\\
Rückwärtsinduktion & vom Blatt zur Wurzel (7)\\
Sattelpunkt & Zeilenmin und Spaltenmax (1.4)\\
Shapley-Wert & mittlerer Beitrag (14.5)\\
Stabile Zuordnung & kein blockierendes Paar (13.3)\\
Teilspielperfekt & Nash in jedem Teilspiel (7.3)\\
Tit for Tat & spiegelt den letzten Zug (8)\\
Trägermenge & benutzte Strategien (3)\\
Volkssatz & fast alles ist Gleichgewicht (8.3)\\
Wert des Spiels & gemeinsamer Minimax-Wert (2.2)\\
Zweitpreisauktion & Sieger zahlt das zweite Gebot (13)\\
\bottomrule
\end{longtable}}
\end{center}
\end{merksatz}
